% Evaluation Section for AetherGuard SCI Paper
% Evidence synchronized with software validation on 2026-07-27.

\section{Evaluation}
\label{sec:evaluation}

\subsection{Evidence Scope}
\label{subsec:evidence-scope}

This evaluation reports only measurements reproducible from the repository and the dated machine-readable validation record. It does not claim a production deployment, a physical-hardware experiment, a live large-language-model comparison, or an external IDS/WAF baseline. The evaluated source baseline is commit \texttt{32b39ff25a909a3a32268fdff1dbe208fffdd4f6}, with an uncommitted correction to the dataset burst-latency denominator. The correction divides accumulated latency by the number of requests actually evaluated before an early burst detection, rather than by the planned burst size.

The test host used Linux 7.0.0-28-generic on x86\_64, an Intel Core i5-9300H with eight logical CPUs, and 7.62~GB of memory. The toolchain comprised Node.js 22.22.1 and Rust 1.97.1. Containers were not used. Backend automated tests ran in test mode with \texttt{DISABLE\_REDIS=true}; consequently, those tests do not constitute live Redis, PostgreSQL, S3, identity-provider, notification, or model-service validation.

\subsection{Synthetic Corpus Evaluation}
\label{subsec:synthetic-corpus}

The tracked AetherGuard corpus contains 4,500 synthetic samples: 3,300 labeled attacks and 1,200 labeled normal requests across 31 categories. We evaluated every sample with the local heuristic detection pipeline. Burst-oriented categories such as brute force, DDoS, and port scanning were evaluated as request sequences. The test intentionally enforces no quality threshold; it verifies complete evaluation and emits a machine-readable report.

\begin{table}[htbp]
\centering
\caption{Tracked Synthetic Corpus Results}
\label{tab:synthetic-corpus-results}
\begin{tabular}{|l|r|}
\hline
\textbf{Metric} & \textbf{Measured value} \\
\hline
Samples & 4,500 \\
True positives & 2,111 \\
False negatives & 1,189 \\
True negatives & 1,200 \\
False positives & 0 \\
Attack detection rate & 63.97\% \\
Synthetic-normal false-positive rate & 0.00\% \\
Median evaluation latency & 0.546 ms \\
95th-percentile evaluation latency & 1.557 ms \\
\hline
\end{tabular}
\end{table}

The aggregate result is heterogeneous across categories. Command injection, brute force, DDoS, port scanning, file upload, HTTP smuggling, Log4j injection, and Text4Shell reached 100\% on this corpus, while XSS reached 98\%. Lower results included SQL injection at 64\%, path traversal at 70\%, SSRF at 20\%, JWT attacks at 24\%, encoding bypass at 13\%, type confusion at 11\%, and Unicode homoglyph attacks at 6\%. These measurements show that the corpus does not support the historical 92.2\% aggregate detection claim.

The 0\% false-positive rate applies only to the 1,200 synthetic normal samples. It is not an estimate of field false-positive rates. No Suricata, ModSecurity, Snort, Fail2ban, iptables, CICIDS, or OWASP Benchmark baseline was run on the host, so this paper makes no comparative accuracy or statistical-significance claim from this experiment.

\subsection{Automated Software Validation}
\label{subsec:automated-validation}

The full backend Jest coverage run completed 248 test suites and 3,291 tests with no failures. Measured coverage was 65.90\% statements, 51.41\% branches, 65.90\% functions, and 66.21\% lines. The frontend type check, lint, eight internationalization contracts, credential-hygiene check, and production build passed. The production build emitted a non-blocking warning for a JavaScript chunk larger than 500~kB.

The Rust fastpath passed 17 unit tests, Clippy with warnings denied, and a release build. The independent rescue component passed 51 unit tests, Clippy, a release build, and 12 fail-closed software drills. Gateway simulator scenarios for invalid signatures, monotonic-counter rollback, stale heartbeats, and emergency-stop latching also passed. Simulator and rescue results are software evidence only and do not validate a physical interlock or secure element.

Deployment regression tests covered successful PM2 activation, rollback after readiness failure, rejection of a mismatched source commit, rejection of symbolic links in a release archive, and offline environment backfill. Repository guards also confirmed that 109 frontend API calls resolve to 170 registered backend routes.

\subsection{Local Throughput Measurements}
\label{subsec:throughput}

We built the Rust fastpath in release mode and bound it to the loopback interface. Autocannon and the server shared the same eight logical CPUs. The fixed request used one tenant and one key, exercising the contended fixed-key rate-check path. Each recorded round used 200 connections, HTTP/1.1 pipeline depth 10, and a 20-second duration after warm-up.

\begin{table}[htbp]
\centering
\caption{Rust Fixed-Key Loopback Rate-Check}
\label{tab:rust-fixed-key-throughput}
\begin{tabular}{|r|r|r|r|r|}
\hline
\textbf{Round} & \textbf{QPS} & \textbf{Mean (ms)} & \textbf{P99 (ms)} & \textbf{Errors} \\
\hline
1 & 61,552.8 & 32.47 & 78 & 0 \\
2 & 59,200.8 & 33.73 & 78 & 0 \\
3 & 56,581.6 & 35.25 & 79 & 0 \\
\hline
\end{tabular}
\end{table}

The three-round mean was 59,111.73 QPS. A short concurrency sweep peaked at 75,078 QPS with 200 connections and pipeline depth 20, with p99 latency of 108~ms and no errors, timeouts, or non-2xx responses. Increasing concurrency to 400--800 connections increased latency without reaching 100,000 QPS. Earlier interactive testing had observed three values above 100,000 QPS, but its complete parameters were not versioned and the current fixed-key reproduction did not reproduce them. We therefore do not claim stable 100,000-QPS capacity.

For contrast, the running single-process Node.js development server, with request logging enabled, processed 3,119.94, 3,242.47, and 3,202.87 QPS over three 15-second loopback rounds on its public health middleware path. This is a development-profile observation rather than a production capacity benchmark.

\subsection{Browser Acceptance Results}
\label{subsec:browser-results}

The Chromium release subset contained ten runnable tests; eight passed and two failed. One additional test was conditionally skipped because it requires a pitch-enabled frontend build. The first failure exposed a route-contract mismatch: a pitch-disabled protected route redirected to \texttt{/login} when the test expected \texttt{/dashboard}. The second occurred after successful authentication, when the walkthrough observed HTTP 503 responses from \texttt{/api/stats/badges} and \texttt{/api/auth/me}; public liveness and readiness endpoints remained HTTP 200. Firefox, WebKit, Edge, ARM64, Windows, and macOS were not tested.

\subsection{Dependency Audit}
\label{subsec:dependency-audit}

Production-dependency audits reported zero vulnerabilities for the root, backend, E2E, installer, and simulator packages. The frontend audit reported two moderate React Router advisories and no high or critical advisories. The dependency graph was not automatically modified during validation.

\subsection{Limitations and Required Follow-Up}
\label{subsec:limitations}

The current evidence establishes repository-level software behavior, not field effectiveness. Required follow-up includes: live PostgreSQL and Redis validation; S3-compatible Object Lock enforcement; live OIDC, SMTP, Slack, and webhook integrations; real local and cloud model inference; an external OWASP Benchmark corpus; independently configured IDS/WAF baselines; multi-node and cross-platform tests; isolated agent CPU and memory measurement; production TLS and cross-host load generation; long-duration soak testing; and physical hardware-in-the-loop experiments.

Accordingly, this evaluation does not support the prior 96.2\% detection accuracy, 2.3\% false-positive rate, 1,200-server production deployment, 99.97\% uptime, or multi-model ablation claims. Those quantities require new, versioned raw evidence before publication.
